% Pattern: proof/attach/gap. A proof that comes after its statement with one
% paragraph of text between them, which remarks on the statement. A second
% proof comes after a statement with two paragraphs of more than seven lines
% between them, and a third
% after a statement and a section head; both stay orphans.
% Seen in the gold pilot (set E).
% want lemma-1 statement/lemma 1 status=here labels=key,eq:key
% want proof-1 proof - proves=lemma-1
% want theorem-1 statement/theorem 2 status=none
% want proof-2 proof - proves=
% want lemma-2 statement/lemma 3 status=none
% want proof-3 proof - proves=
\documentclass{article}
\newtheorem{theorem}{Theorem}
\newtheorem{lemma}[theorem]{Lemma}
\begin{document}
\begin{lemma}\label{key}
An identity
\begin{equation}\label{eq:key}
a = b.
\end{equation}
\end{lemma}
The identity \eqref{eq:key} is used throughout, and
\begin{equation*}
c = d
\end{equation*}
is a special case.
\begin{proof}
A step.
\end{proof}
\begin{theorem}
A statement.
\end{theorem}
A first paragraph of text,
which runs on
over three lines.

A second paragraph of text,
which runs on
over three lines too.

\begin{proof}
Another step.
\end{proof}
\begin{lemma}
A helper.
\end{lemma}
\section{More}
\begin{proof}
A third step.
\end{proof}
\end{document}
